% ======================================================================
\section{Known limitations and open issues}
\label{sec:limitations}
% ======================================================================

This section lists the limitations that bear on the path from the field to
the molecules, with the direction of the bias each introduces and its
measured size where one exists. The first group is recorded in the project's
own logs. The second group was found while reading the code for this document
and is not, as far as could be found, recorded in \src{DEVLOG.md},
\src{DEVLOG-H2.md} or \src{doc/\allowbreak{}pdr-status.md}; its sizes have not been
measured.

\subsection{Recorded in the project's logs}

\begin{enumerate}
\item \textbf{Per-cell Monte Carlo noise in the photorates} (\dlog{L-26},
  \dlog{D-405}, \dlog{H-012}). Every photorate reads its own cell's field, whose
  per-cell scatter reaches 14--22\% beyond $\AV\sim2$ at the default packet
  count; the level model's per-line field is noisier than the band integral.
  Besides scatter, the noise biases cell means through the non-linear rates
  (estimated $+5\%$ at 22\% scatter) and, at the carbon turn, decides which of
  two states a cell is in (\dlog{D-405}). An estimator accumulated over
  iterations is registered for future development (\dlog{D-411}).
  Section~\ref{sec:mcnoise} gives the numbers.
\item \textbf{Packet convergence of carbon and CO} (\dlog{D-408}, \dlog{H-032}).
  At the default $6\times10^6$ packets the \Htwo{} lines and the H/\Htwo{}
  crossing are converged, but the carbon turn is placed too shallow (C$^+=$C
  moves by 0.5\,mag on \texttt{hiipdr} up to $3.84\times10^8$ packets, towards
  Cloudy's), the deep gas is too cold, [C\,\textsc{i}] needs 16--64 times the
  packets and CO is still moving at 64 times. The stop test does not see this.
\item \textbf{The twelve-direction quadrature} (\dlog{L-53}, \dlog{D-410}):
  cells that the grid's symmetry makes equal differ in $n(\mathrm{H_2})$ by up
  to about 11\% following the ring axis of the twelve directions, independently
  of the packet count; 48 directions remove the pattern at 1.48 times the run
  time, moving totals by under 1\%. The twelve were kept by decision, since no
  line map carries the pattern (CO's abundance and the \Htwo{} density do).
\item \textbf{Damping width of the \Htwo{} lines} (\dlog{H-008}, \dlog{H-026}).
  The code damps each line with the upper level's total radiative decay rate
  (to the bound levels and the continuum of $X$), which is the physical natural
  width; damping each line
  with its own $A$, as Cloudy does, narrows the wings, strengthens the
  shielding and moves the H/\Htwo{} crossing 0.21\,mag shallower on \texttt{hiipdr}
  (from $\AV=1.396$ to 1.188, Cloudy's being 1.280) and 0.10\,mag shallower on
  \texttt{pn150k}. Between $\AV=1.05$ and 1.2 on \texttt{hiipdr} the model's
  photodissociation rate per molecule is 5--7 times Cloudy's. On \texttt{hiipdr}
  the width accounts for $+0.21$\,mag of the difference between the two codes'
  crossings and everything else for $-0.09$\,mag (\dlog{H-026}); it is the
  largest single identified term and should accompany any comparison with
  Cloudy, which uses the total decay rate for every species but \Htwo{}
  (\dlog{H-008}).
\item \textbf{No H\,\textsc{i} Lyman-series shading} of the \Htwo{} band
  (\src{DEVLOG-H2.md}, \dlog{H-006}, ``What is not yet known''). Neither the MC transfer
  nor the level model contains it, so pumping and dissociation are
  overestimated where a large H\,\textsc{i} column lies in front of the molecular
  gas, and the transition is placed too deep. The hydrogen column between the
  ionisation front and the H/\Htwo{} crossing is about $1.8\times10^{21}$\,cm$^{-2}$ on
  \texttt{hiipdr} ($\AV$ 0.48 to 1.40 at $1.93\times10^{21}$\,cm$^{-2}$\,mag$^{-1}$,
  Table~\ref{tab:dustfix}; \dlog{D-260}, \dlog{H-027}) and at most
  $1.1\times10^{21}$\,cm$^{-2}$ on \texttt{pn150k} (crossing at $\AV=0.39$,
  $2.79\times10^{21}$\,cm$^{-2}$\,mag$^{-1}$), part of it molecular. The damping wings
  of Ly$\beta$, Ly$\gamma$ and Ly$\delta$ alone cover about 4, 1.4 and 0.7\,\AA{} of the
  912--1110\,\AA{} band at $N(\mathrm{H\,I})=10^{21}$\,cm$^{-2}$, growing as $N^{1/2}$
  (an estimate, Section~\ref{sec:overlap}). Unmeasured.
\item \textbf{No cross-pumping and no overlap in trapping} (\dlog{H-006}): a
  photon absorbed in another line's wing is lost rather than pumping that
  line, and the escape of each line ignores the others.
\item \textbf{Electrons are not colliders} of the \Htwo{} levels, and the level
  model has no destruction other than Solomon (\src{h2\_mod.f90} header); this
  matters in the partly ionised gas, where \Htwo{} is already 14--230 times
  Cloudy's for chemical reasons (\dlog{L-43}).
\item \textbf{C\,\textsc{i} photoionisation without \Htwo-line shielding}
  (\src{DEVLOG.md}, item 8 of the audit list near line 28742): MOCASSIN's rate
  replaces \rtn{CIPDRATE}, and MOCASSIN does not transport LW lines, so the
  rate is overestimated behind the H/\Htwo{} transition; estimated at a few
  tens of per cent at $\AV\sim3$, with the opposite sign to the deep \Cp{}
  deficit (\dlog{L-3}).
\item \textbf{Chlorine is absent from the network}; in Cloudy it carries 25--50\%
  of the \Htwo{} sink from $\AV=0.26$ on in the PN model, and Cloudy without
  chlorine turns molecular 0.042\,mag earlier (\dlog{L-46}, \dlog{H-026}). Its
  absence moves the model's crossing shallower, by about 0.04\,mag on
  \texttt{pn150k}, partly offsetting there the 0.10\,mag of item~4 (\dlog{L-46}).
\item \textbf{Silicate formation efficiency}: the per-material Cazaux \& Tielens
  surface energies of 3D-PDR give 0.50--0.91 of Cloudy's coefficient on
  \texttt{hiipdr} (\dlog{L-33}, \dlog{H-022}); reported, not tuned.
\item \textbf{The mean grain temperature can switch formation off}
  (\dlog{H-035}): the number-weighted mean over the whole size grid counts the
  bins outside a species' size range at zero temperature; on a grid wider
  than the grains present this put $T_d$ at $10^{-6}$\,K and no \Htwo{} formed.
  Fixed after \texttt{8fc1512}: \dlog{D-414} takes the mean over the bins that hold
  grains (Section~\ref{sec:scope}).
\item \textbf{Diffuse-site weights} recover 52--130\% of the MC's diffuse field
  (\dlog{L-4}); the magnitude is the MC's, the directional distribution, and so
  the columns assigned to the diffuse light, is approximate
  (Section~\ref{sec:sites}).
\item \textbf{Unresolved fronts} (\dlog{L-37}, \dlog{D-088}): on uniform grids
  the ionisation front and the H/\Htwo{} transition are thinner than a cell; the
  front cells, which carry the warm partly ionised gas through their whole
  volume, make 92\% of \texttt{hiipdr}'s [O\,\textsc{i}] 63\,\textmu m, 9.3--9.5 times
  Cloudy's at $31^3$ and 5.5 times at $61^3$. Front-sensitive lines need a
  resolution check; a sub-cell estimate of [O\,\textsc{i}] 6300\,\AA{} is available
  offline (\kw{frontLines}).
\item \textbf{The gas re-emits from the hydrogen and helium spectrum whatever
  absorbed} (\dlog{L-51}). A packet absorbed by neutral carbon in the PDR is
  re-emitted, some of the time, as a hydrogen- or helium-ionising photon:
  1.9\% of \texttt{hiipdr}'s cells per iteration at $\AV=2.15$--2.43 are
  over-ionised, and their chemistry, solved to steady state on that iteration's
  ionisation, can lose more than half its \Htwo{} (202 such drops in six
  iterations at the default packet count, 23 at four times, none at sixteen
  times).
\item \textbf{No PAHs or very small grains} (\dlog{L-7}): the photoelectric
  heating lacks the component that dominates in real PDRs, so the gas is colder
  than in PDR models with that component at the same field; common-mode with
  the Cloudy references.
\item \textbf{Gas--grain collisional coupling omits \Htwo} (\dlog{L-9}): 68
  times weaker than Cloudy's dust cooling at $\AV=0.65$ on \texttt{pn150k},
  negligible at $n_{\mathrm H}=10^3$\,cm$^{-3}$ and decisive at high density.
\item \textbf{The \Htwo{} line emission is computed optically thin and without
  dust extinction} in \src{lineFluxH2.out} (\src{h2Lines\_mod.f90} header).
\end{enumerate}

\subsection{Found while preparing this document}

\begin{enumerate}\setcounter{enumi}{17}
\item \textbf{Formation flux in the level model is tied to the Solomon flux.}
  \rtn{solveCell} calls \rtn{h2Solve} without a sink or a source
  (\src{pdr\_mod.f90:\allowbreak{}5904}), so the rate at which new molecules enter the
  levels in the formation distribution $\phi$ equals the Solomon dissociation
  rate. The chemistry's actual formation rate, $R_{\mathrm f}n_{\mathrm H}n(\mathrm H)$,
  balances all destruction (cosmic rays, ion--molecule reactions, Solomon) in
  steady state. Where the Solomon rate has been shielded away, formation
  pumping is therefore underestimated. With the network's \Htwo\,+\,CRP channels
  alone the other destruction is $0.99\,\zeta$, $10^{-16}$\,s$^{-1}$ at the default
  \kw{pdrZeta}, so this happens where the Solomon rate is shielded below a few
  $10^{-16}$\,s$^{-1}$ (Section~\ref{sec:balance}). The energy budget is split the same
  way: the internal energy that newly formed molecules bring to the levels
  (and that the collisional term returns to the gas) follows the Solomon flux,
  while the kinetic part of the formation energy (\texttt{heat(5)}) follows the
  chemistry's formation rate. Both numbers needed for a consistent closure are
  already available in \rtn{solveCell} ($w_{\mathrm{sil}}+w_{\mathrm{gra}}$,
  $n(\mathrm H)$, $n(\mathrm{H_2})$): passing the non-Solomon destruction as
  \rtn{sink}, or the chemistry's formation flux per molecule as \rtn{src}, would
  make the two agree. A per-cell print of formation flux over Solomon flux
  would size the effect first.
\item \textbf{The final output mixes two Solomon rates.} The last
  \rtn{h2LevelModel(grid,12)} (\src{iteration\_mod.f90:\allowbreak{}3254, 3283}) runs after
  the last chemistry, which is not re-solved. The written level populations and
  $k_{\mathrm{Sol}}$ are those of up to twelve sweeps; the written
  $n(\mathrm{H_2})$ was obtained with the previous one-sweep rate. The size of
  the difference is printed by the call itself (``Solomon rate against the
  last call's'').
\item \textbf{The \Htwo{} collisional heating term is frozen and two
  iterations old in the temperature search.} $\mathcal{H}_{\mathrm{coll}}$ changes
  sign and depends steeply on $T$ (it is \Htwo's line cooling where the gas is
  warmer than the levels), but it enters \rtn{pdrHeatXtra}, which MOCASSIN holds
  constant in \rtn{iterateT}, and the value used to find $T_e^{(n)}$ was
  computed by the level model at $T_e^{(n-2)}$ (Section~\ref{sec:lags}).
  \dlog{H-030} associates this term, together with the damping, with the two
  derailed iterations at the start of the chemistry. Tabulating it against
  $T$, as the CO cooling is (\dlog{D-215}), would remove the freeze.
\item \textbf{Generic photoreactions keep the Draine spectral shape}
  (Section~\ref{sec:generic}). \dlog{L-24} corrected the LW-window reactions;
  the other 24 photoreactions still multiply the 6--13.6\,eV band strength by
  coefficients computed for the Draine shape, although the MC field's shape
  changes strongly with depth (its LW share falls from 0.30 to 0.008 on
  \texttt{hiipdr}, \dlog{D-349}). Reactions whose thresholds lie in the window
  are overestimated at depth, those driven by the 6--11\,eV part underestimated.
  Folding each reaction's cross section, or a few sub-bands, over the MC field
  would remove the assumption.
\item \textbf{The level model's \Htwo{} transmission is not shared.} The
  per-line \Htwo{} shielding computed in \rtn{solveCell} is exactly what is
  missing from MOCASSIN's C\,\textsc{i} rate (item~8) and what CO's
  van Dishoeck \& Black factor approximates for \Htwo{} screening; neither uses
  it.
\item \textbf{CO shielding and the chemistry's own \Htwo{} diagnostic} exclude
  the target cell's own column (\dlog{D-088}); only the level model averages
  through it. On \texttt{hiipdr}'s grid half a cell can hold $10^{16}$\,cm$^{-2}$ of
  CO, so CO is under-shielded in the cells of its own transition and the
  transition is placed up to about half a cell too deep
  (Section~\ref{sec:avradial}).
\item \textbf{Several sources are combined by geometry alone.} The effective
  $\AV$, the source-ward columns, the in-cell chord (\rtn{computeAVradial},
  \src{pdr\_mod.f90:\allowbreak{}1986}) and the level model's source-ward level columns
  (\rtn{starColumns}, \src{:\allowbreak{}5651}) weight the sources by $r_s^{-2}$, without
  their luminosities or their attenuation, and apply the mean column to the
  summed direct field through non-linear shielding functions. Exact for one
  source, which is every fixture; untested for several. Weighting each source
  by its own attenuated direct field, $L_{\mathrm{FUV},s}r_s^{-2}e^{-\RFUV\tau_{V,s}}$,
  and evaluating the shielding per source would remove it.
\item \textbf{$\RFUV$ is weighted by the emitted spectrum of source~1}
  (\src{pdr\_mod.f90:\allowbreak{}1381--1390}). It overstates the band extinction of the
  local, reddened field at depth, and so the optical depth between the
  scattering sites and the target; this suppresses distant sites increasingly
  with depth and is a candidate for the fall with depth of the formal
  solution's recovery in \dlog{L-4}.
\item \textbf{The grain temperature is a number-weighted mean applied to an
  area-weighted rate} (Section~\ref{sec:Td}). Besides the zero-bin trap of
  item~11, the mean favours the smallest grains and one efficiency
  $\epsilon(\langle T_d\rangle)$ replaces $\langle\epsilon(T_d)\rangle$ over the sizes. On
  silicate at 20--50\,K, where $\epsilon$ falls by a factor of 3.6
  (Section~\ref{sec:Td}), the formation rate depends on how the temperatures
  of the sizes are spread. Evaluating the rate per size with $\pi a^2$ weights
  would remove both; a temperature per grain material is registered for
  future development as \dlog{D-417} (\dlog{L-54}), after the commit described
  here.
\item \textbf{No interstellar field at the outer boundary}
  (Section~\ref{sec:assumptions}, item~6). On \texttt{hiipdr} the source's field
  is below $\chi=1$ in the outer 40\% of the radius, so the outer PDR is more
  molecular and colder than the same cloud in the interstellar field.
\item \textbf{The coolant line escape depends weakly on the sweep order.} The
  optical depths of the coolant lines read the other cells' level populations
  as stored on the rank (\src{pdr\_mod.f90:\allowbreak{}4315--4316}), which cells solved
  earlier on the same rank have already updated
  (\src{:\allowbreak{}4440--4442}); the shielding columns are snapshots and do not have
  this dependence.
\item \textbf{The seed double-attenuates.} \rtn{seedMolecules} multiplies the
  local, already attenuated $\chi$ by \rtn{SCATTER}$(\AV)$
  (\src{pdr\_mod.f90:\allowbreak{}2937--2938}; Section~\ref{sec:seed}). It affects only the
  starting point of the first PDR iteration.
\item \textbf{The plane-parallel branch includes the near half of the cell in
  its columns} (\src{pdr\_mod.f90:\allowbreak{}2492--2507}), against \dlog{D-088}'s rule in
  the point-source branch. The level model does not run in that geometry.
\item \textbf{Collisional ionisation of hydrogen lacks the electron density.}
  MOCASSIN adds $2.75\times10^{-16}T^{3/2}(157\,894/T+2)e^{-157\,894/T}$, a rate
  coefficient in cm$^3$\,s$^{-1}$, to the photoionisation rate in s$^{-1}$ without the
  factor $n_e$, in its ionisation balance (\src{update\_mod.f90:\allowbreak{}2769--2770})
  and in the rate handed to the chemistry (\src{:\allowbreak{}2859}). Collisional
  ionisation is therefore misscaled by a factor $n_e$ (underestimated wherever
  $n_e>1$\,cm$^{-3}$); with the factor, the term would reach 10\% of a
  photoionisation rate of $10^{-8}$\,s$^{-1}$ at about $1.5\times10^4$\,K for
  $n_e=10^4$\,cm$^{-3}$, so hot dense ionised gas is affected and the PDR is not.
  Inherited from MOCASSIN (Section~\ref{sec:linked}).
\item \textbf{Two reactions of the rate file are inactive.} The grain-assisted
  recombinations of He$^+$ and C$^+$ (reactions 330 and 331) have no term in
  \src{pdr/\allowbreak{}odes\_reduced.c}; their coefficients are computed but unused
  (Section~\ref{sec:network}).
\item \textbf{H\,+\,\Htwo{} rates below 100\,K} are those of 100\,K, the first
  temperature of the \citet{Lique2015} table, which overestimates the reactive
  ortho--para conversion by H in cold gas, where \Hp{} exchange usually
  dominates it (Section~\ref{sec:coll}).
\item \textbf{The \Htwo{} line profile is normalised to its own centre}
  ($\tau_0H(a,x)/H(a,0)$ with Doppler-core $\tau_0$), overstating the Voigt optical
  depth by $1/H(a,0)$, at most 3\% for these data (Section~\ref{sec:shield}).
\item \textbf{Stale or inaccurate comments.} \src{constants\_mod.f90:\allowbreak{}18--20} labels
  the FUV lower edge 2400\,\AA{} (the value is 6.0\,eV, 2068\,\AA); the comment in
  \rtn{addOpacity} (\src{ionization\_mod.f90:\allowbreak{}444--453}) describes a power-law
  \Htwo{} cross section, whereas \rtn{yanH2xSec} is the piecewise fit of
  \citet{Yan1998}; and the comment of \rtn{draineWindowShare}
  (\src{pdr\_mod.f90:\allowbreak{}2285--2286}), \dlog{D-349} and \dlog{L-24} give the Draine
  field's LW share of the continuous band as 0.1516 (0.152), whereas
  Equation~(\ref{eq:draine}) gives 0.148. The code uses the binned value
  (0.151--0.164 on the fixtures), consistently with the MC sums, so no result
  is affected. The comment of \rtn{h2FormDist} (\src{h2\_mod.f90:\allowbreak{}578--579})
  gives the vibrational distribution as ``their eqs 4-5'' about
  $E_{vm}=X_{\mathrm{vib}}D$ and the rotational one as ``their eqs 6 and 7'', whereas
  in \citet{TakahashiUehara2001} the vibrational distribution is eqs.~5--7 (with
  $E_{vm}=X_{\mathrm{vib}}D+E_{\mathrm{off}}$) and the rotational models are eqs.~12
  and 13 (the code is right); and \src{update\_mod.f90:\allowbreak{}2621} attributes the
  collisional-ionisation formula to \citet{DrakeUlrich1980}, which does not
  contain it (Section~\ref{sec:linked}).
\end{enumerate}
